\documentclass[12pt, letterpaper]{article}

% ── Packages ──────────────────────────────────────────────────────────────────
\usepackage[margin=1in, headheight=15pt]{geometry}
\usepackage{amsmath, amssymb}
\usepackage{fontenc}
\usepackage{inputenc}
\usepackage{microtype}
\usepackage{setspace}
\usepackage{parskip}
\usepackage{titlesec}
\usepackage{titling}
\usepackage{fancyhdr}
\usepackage{hyperref}
\usepackage{xcolor}
\usepackage{enumitem}
\usepackage{graphicx}
\usepackage{caption}
\usepackage{subcaption}
\usepackage{abstract}
\usepackage{float}

\graphicspath{{figures/}}

% ── Colors ────────────────────────────────────────────────────────────────────
\definecolor{headingcolor}{RGB}{30, 30, 30}
\definecolor{rulecolor}{RGB}{120, 120, 120}
\definecolor{darkgray}{RGB}{80, 80, 80}
\definecolor{linkblue}{RGB}{30, 90, 160}

% ── Hyperref setup ────────────────────────────────────────────────────────────
\hypersetup{
    colorlinks=true,
    linkcolor=linkblue,
    citecolor=linkblue,
    urlcolor=linkblue,
    pdftitle={Neuro-Inspired Computing on the SoC FPAA: Spring 2026 Undergraduate Research Report},
    pdfauthor={Ethan Mazor}
}

% ── Typography ────────────────────────────────────────────────────────────────
\onehalfspacing
\setlength{\parindent}{0pt}
\setlength{\parskip}{0.75em}

% ── Section formatting ────────────────────────────────────────────────────────
\titleformat{\section}
  {\Large\bfseries\color{headingcolor}}
  {\thesection.}{0.75em}{}
  [\color{rulecolor}\vspace{0.2em}\hrule height 0.6pt\vspace{0.2em}]

\titleformat{\subsection}
  {\large\bfseries\color{headingcolor}}
  {\thesubsection}{0.75em}{}

\titlespacing{\section}{0pt}{1.5em}{0.75em}
\titlespacing{\subsection}{0pt}{1em}{0.4em}

% ── Header / Footer ───────────────────────────────────────────────────────────
\pagestyle{fancy}
\fancyhf{}
\fancyhead[L]{\small\color{darkgray}\textit{Neuro-Inspired Computing on the SoC FPAA}}
\fancyhead[R]{\small\color{darkgray}\textit{Ethan Mazor}}
\fancyfoot[C]{\small\color{darkgray}\thepage}
\renewcommand{\headrulewidth}{0.4pt}
\renewcommand{\footrulewidth}{0pt}
\fancypagestyle{plain}{%
  \fancyhf{}
  \fancyfoot[C]{\small\color{darkgray}\thepage}
  \renewcommand{\headrulewidth}{0pt}
}

% ── Abstract formatting ───────────────────────────────────────────────────────
\renewcommand{\abstractnamefont}{\large\bfseries\color{headingcolor}}
\renewcommand{\abstracttextfont}{\normalsize}
\setlength{\absleftindent}{0.5in}
\setlength{\absrightindent}{0.5in}

% ── List / caption formatting ─────────────────────────────────────────────────
\setlist[itemize]{leftmargin=1.5em, itemsep=0.3em, topsep=0.3em}
\captionsetup{font=small, labelfont=bf, skip=5pt}

% ─────────────────────────────────────────────────────────────────────────────
\begin{document}

% ── Title Page ────────────────────────────────────────────────────────────────
\begin{titlepage}
  \centering
  \vspace*{2cm}

  {\color{headingcolor}\rule{\linewidth}{2pt}}
  \vspace{1em}

  {\Huge\bfseries\color{headingcolor}
    Neuro-Inspired Computing\\[0.4em]
    on the SoC FPAA
  \par}
  \vspace{0.6em}
  {\Large\color{darkgray} Spring 2026 Undergraduate Research Report\par}
  \vspace{0.5em}
  {\color{rulecolor}\rule{\linewidth}{1.5pt}}

  \vspace{2.5cm}

  {\large Ethan Mazor\par}
  \vspace{0.4em}
  {\normalsize\color{darkgray}
    Georgia Tech-Europe\\
    School of Electrical and Computer Engineering\\[0.6em]
    Advisor: Dr.\ Alexandre Locquet\\[0.6em]
  \par}

  \vfill
\end{titlepage}

\thispagestyle{plain}

% ── Abstract ─────────────────────────────────────────────────────────────────
\begin{abstract}
This report documents progress made during the Spring 2026 semester of undergraduate research in the field of neuro-inspired analog computing. Work centered on implementing Hopfield networks on a Field Programmable Analog Array (FPAA) and using it to solve the Max-Cut problem. The primary experimental focus was investigating ways of escaping local energy minima via externally applied "kicks." The report covers a high-level overview of the FPAA platform, Hopfield networks and their implementation on the FPAA, the local vs.\ global minimum problem, experimental exploration of kick signal parameters, and lessons learned over the semester.
\end{abstract}

% ─────────────────────────────────────────────────────────────────────────────
\section{The Field Programmable Analog Array (FPAA)}
\label{sec:fpaa}

The SoC FPAA is a reconfigurable analog computing chip that is analogous to its more well-known counterpart, the FPGA. It can be thought of as a universal analog breadboard: rather than physically wiring components, a user specifies a circuit in software and the chip programs itself. The core building blocks are Computational Analog Blocks (CABs) which contain several analog components such as OTAs, FG OTAs, etc. The specific boards available were the 3.0 A boards, chip numbers 11 and 12.

The design and programming workflow runs through the \textbf{CADSP} software where circuits are assembled from functional blocks. Xcos compiles the design into a hardware configuration that then gets programmed onto the board's chip. Physical interfacing and measurements were made using the Digilent \textbf{Analog Discovery} versions 2 or 3, which served as both a waveform generator and an oscilloscope.

At the start of the semester, simple circuits were built to learn the toolchain before moving on to Hopfield networks. Figure~\ref{fig:lpf_early} shows the first Xcos design (a low-pass filter) alongside a DAC-to-ADC calibration sweep. The LPF was then verified on the physical
board; Figure~\ref{fig:lpf_meas} shows the filtered square wave output and the Bode magnitude plot from a frequency sweep.

\begin{figure}[H]
  \centering
  \begin{subfigure}[b]{0.52\textwidth}
    \includegraphics[width=\textwidth]{lpf_block}
    \caption{First Xcos design: a two-OTA low-pass filter.}
  \end{subfigure}
  \hfill
  \begin{subfigure}[b]{0.44\textwidth}
    \includegraphics[width=\textwidth]{dac_adc}
    \caption{DAC-to-ADC sweep showing analog readback variability.}
  \end{subfigure}
  \caption{Early FPAA exercises used to learn the software toolchain and hardware behavior.}
  \label{fig:lpf_early}
\end{figure}

\begin{figure}[H]
  \centering
  \begin{subfigure}[b]{0.50\textwidth}
    \includegraphics[width=\textwidth]{lpf_squarewave}
    \caption{Filtered square wave output measured on-board.}
  \end{subfigure}
  \hfill
  \begin{subfigure}[b]{0.46\textwidth}
    \includegraphics[width=\textwidth]{bode}
    \caption{Bode magnitude plot, 10--500\,Hz.}
  \end{subfigure}
  \caption{Physical LPF measurements taken using the Analog Discovery.}
  \label{fig:lpf_meas}
\end{figure}

% ─────────────────────────────────────────────────────────────────────────────
\section{Hopfield Networks on the FPAA}
\label{sec:hopfield}

A Hopfield network is a fully-connected recurrent neural network which is designed to
solve problems by mirroring how physical systems naturally settle into states of minimum
energy. The network's state and energy function are:

\begin{equation}
  S_i = \begin{cases} 1 & \text{if } \sum_{j} W_{ij} S_j \geq T \\ 0 & \text{otherwise}
        \end{cases}
  \qquad\qquad
  E(\mathbf{x}) = -\tfrac{1}{2}\,\mathbf{x}^{\top} W \mathbf{x}
                  - \boldsymbol{\theta}^{\top} \mathbf{x}
  \label{eq:hopfield}
\end{equation}

In the above energy function, $W$ represents the weight matrix, which functions as the
system's memory by defining the connections between every pair of neurons in the network.
The vector $\mathbf{x}$ denotes the current state of all neurons in the network (high or
low). Finally, the $\theta$ term represents an external bias, acting as a constant signal
that shifts the energy landscape to favor certain states over others.

\subsection{The Max-Cut Problem}

The optimization problem targeted in this research was the Max-Cut problem. Given an
undirected graph, the goal is to partition the nodes into two sets such that the number
of edges crossing between the two sets (the ``cut'') is maximized. Max-Cut is NP-hard, and
maps naturally onto a Hopfield network by assigning one neuron per graph node: a neuron's
binary state (0 or 1) indicates which of the two sets that node belongs to. By constructing
the weight matrix such that connected nodes are penalized for sharing the same state, the
network's stable states correspond to high-quality cuts of the graph.

\begin{figure}[H]
  \centering
  \includegraphics[width=0.55\textwidth]{maxcut}
  \caption{Max-Cut on a fully-connected 4-node graph. Nodes are split into Set A
           (\{N1, N3\}) and Set B (\{N2, N4\}), partitioned by the dashed line. Solid
           green edges cross the partition (cut); dashed gray edges remain within a set
           (uncut). The optimization goal is to maximize the number of cut edges.}
  \label{fig:maxcut}
\end{figure}

\subsection{Sample Weight Matrix for a 4-Neuron Network}

The weight matrix used to program the FPAA has fixed dimensions of $8 \times 6$ for up
to 4 neurons. Its structure is as follows:

\begin{itemize}
  \item \textbf{Rows} come in pairs $(2k{-}1, 2k)$ for each neuron $N_k$, representing
  the positive ($p$) and negative ($n$) conductance branches of the differential pair
  encoding.

  \item \textbf{Columns W1 and W2} are the bias / forcing inputs, used to drive each
  neuron into a desired initial state. Both columns can also be used to apply external
  forcing signals during computation.

  \item \textbf{Columns 3--6} are the recurrent inputs from neurons $N_1$ through
  $N_4$, defining the connectivity between every pair of neurons.

  \item \textbf{Entries} are scalar current values selected from three ranges.
  Around $1\,\text{pA}$ is treated as effectively zero and indicates no connection;
  values in the $100$--$300\,\text{nA}$ range are used for recurrent weights between
  neurons; and values in the $\mu\text{A}$ range are used in columns W1 and W2 as
  forcing biases. All current values are non-negative; effective sign is encoded by
  the $p$/$n$ row pair. For recurrent weights, a \textbf{negative} differential pair
  ($n > p$) represents an edge between the two neurons (inhibitory), while a
  \textbf{positive} differential pair ($p > n$) represents the absence of an edge
  (excitatory). The total total current across both rows of a single neuron must remain
  under $10\,\mu\text{A}$.
\end{itemize}

A worked example for a 4-neuron star topology (edges N1--N4, N2--N4, N3--N4) is shown
in Figure~\ref{fig:star_with_matrix}. Column W1 is used to set the initial state of the
network --- in this case $\texttt{0000}$, with all neurons forced LOW at
$3\,\mu\text{S}$. Column W2 is reserved for the external kick / forcing signal, which
in this example applies a positive bias of $5\,\mu\text{S}$ to N3 only.

\begin{figure}[H]
  \centering
  \begin{minipage}[c]{0.30\textwidth}
    \centering
    \includegraphics[width=\textwidth]{star_topology}
  \end{minipage}
  \hfill
  \begin{minipage}[c]{0.66\textwidth}
    \centering
    \small
    \setlength{\tabcolsep}{4pt}
    \begin{tabular}{l|cccccc}
           & W1            & W2            & $\to$N1       & $\to$N2       & $\to$N3       & $\to$N4       \\ \hline
    N1$_p$ & 1\,pA         & 1\,pA         & 1\,pA         & 100\,nA       & 100\,nA       & 1\,pA         \\
    N1$_n$ & 3\,$\mu$A     & 1\,pA         & 1\,pA         & 1\,pA         & 1\,pA         & 100\,nA       \\
    N2$_p$ & 1\,pA         & 1\,pA         & 100\,nA       & 1\,pA         & 100\,nA       & 1\,pA         \\
    N2$_n$ & 3\,$\mu$A     & 1\,pA         & 1\,pA         & 1\,pA         & 1\,pA         & 100\,nA       \\
    N3$_p$ & 1\,pA         & 5\,$\mu$A     & 100\,nA       & 100\,nA       & 1\,pA         & 1\,pA         \\
    N3$_n$ & 3\,$\mu$A     & 1\,pA         & 1\,pA         & 1\,pA         & 1\,pA         & 100\,nA       \\
    N4$_p$ & 1\,pA         & 1\,pA         & 1\,pA         & 1\,pA         & 1\,pA         & 1\,pA         \\
    N4$_n$ & 3\,$\mu$A     & 1\,pA         & 100\,nA       & 100\,nA       & 100\,nA       & 1\,pA         \\
    \end{tabular}
  \end{minipage}
  \caption{Left: 4-neuron star topology with edges N1--N4, N2--N4, N3--N4. Right: the
           corresponding $8 \times 6$ FPAA weight matrix. Diagonal entries (e.g.\ N1's
           row pair under column $\to$N1) are set to $1\,\text{pA}$ on both rows,
           enforcing no self-connections.}
  \label{fig:star_with_matrix}
\end{figure}

% ─────────────────────────────────────────────────────────────────────────────
\section{Local vs.\ Global Solutions and Annealing}
\label{sec:localvsglobal}

\subsection{The Local Minimum Problem}

While the energy function in \eqref{eq:hopfield} guarantees convergence, it only guarantees
convergence to a \textit{local} minimum---not necessarily the global optimum. Which minimum the network finds depends heavily on
initial conditions, and a poor starting state can trap the system in a shallow local
minimum far from the optimal solution. This is the central practical challenge of using
Hopfield networks for optimization, and it was the primary experimental focus of the
semester.

\subsection{Analog Annealing via External Kick Signals}

The annealing approach was to apply an external signal through the forcing input columns
of the weight matrix to temporarily perturb the network's energy landscape while stuck in
a local minimum. Once the perturbation is removed, the network reconverges, ideally to a
better solution. Figure~\ref{fig:kick_concept} illustrates this experimentally: the network
is forced into an intended initial state (0010), but the analog hardware settles into a
different local minimum (0101). A kick signal is then applied, and the network reconverges
--- in this case returning to the same state, demonstrating the challenge of reliably
escaping a local minimum.

\begin{figure}[H]
  \centering
  \includegraphics[width=0.95\textwidth]{kick_concept}
  \caption{Demonstration of the kick concept on a 4-neuron Max-Cut network. The intended initial state (0010) diverges from the actual settled state (0101). The external kick (W2) is applied and the network reconverges to 0101 after the perturbation is removed.}
  \label{fig:kick_concept}
\end{figure}

\subsection{Kick Signal Parameters and Experimental Results}

A kick signal is defined by several parameters, each of which influences whether the
perturbation can successfully drive the network out of a local minimum. The parameters
explored during the semester are described below; experimental results are placed
alongside the parameter they primarily vary. All experiments were performed on a
fully-connected 4-neuron network attempting to escape the local solution \texttt{0001}.

\textbf{Signal type.} A square pulse creates a hard step, while a sine pulse provides a
smooth gradual tilt and return. Smoother signals may give the network more time to
reorganize before the perturbation ends; in practice, a square pulse was most commonly
used for the external kick.

\begin{figure}[H]
  \centering
  \includegraphics[width=0.74\textwidth]{kick_sine}
  \caption{Half-sine pulse, peak amplitude $\approx 2.5\,\text{V}$, duration $\approx 30\,\mu$s.
           The network reconverged to the local solution \texttt{0001}.}
  \label{fig:kick_sine}
\end{figure}

\begin{figure}[H]
  \centering
  \includegraphics[width=0.74\textwidth]{kick_chirp}
  \caption{Rectified sine chirp, peak amplitude $\approx 2.5\,\text{V}$, frequency
           increasing over time. The network returned to \texttt{0001}.}
  \label{fig:kick_chirp}
\end{figure}

\textbf{Rise time ($t_r$).} The finite rise time of the signal determines how quickly the
landscape shifts. A very fast rise ($t_r \ll \tau_\text{network}$) may not allow the
network to respond coherently, while a slower rise gives neurons time to begin re-settling
while the perturbation is still active. A very short rise time (relative to convergence
time) was typically used.

\begin{figure}[H]
  \centering
  \includegraphics[width=0.74\textwidth]{kick_tr01}
  \caption{Square pulse, $t_r = 0.1\,\mu$s, peak amplitude $\approx 2.5\,\text{V}$,
           duration $\approx 50\,\mu$s. The network reconverged to \texttt{0001}.}
  \label{fig:kick_fast}
\end{figure}

\begin{figure}[H]
  \centering
  \includegraphics[width=0.74\textwidth]{kick_tr120}
  \caption{Square pulse, $t_r = 120\,\mu$s, peak amplitude $\approx 2.5\,\text{V}$,
           duration $\approx 300\,\mu$s. The network again returned to \texttt{0001}.}
  \label{fig:kick_slow}
\end{figure}

\textbf{Amplitude.} Controls how much the energy landscape is distorted. Too small an
amplitude fails to overcome the local minimum's energy barrier; too large may overwhelm
the recurrent dynamics and drive the network to an arbitrary state. Experimentally, a
$\approx 2.5\,\text{V}$ amplitude signal was found to be most effective for perturbing
the system.

\textbf{Duration / duty cycle.} The total time the perturbation is applied. Too brief a
kick may not allow re-exploration; too long risks convergence to a different local
minimum within the perturbed landscape itself.

\textbf{Target neuron(s).} Rather than perturbing all forcing inputs equally, biasing
specific neurons in the weight matrix tests whether directed perturbations more
effectively steer the network toward a preferred solution. Experimentally, this was found
to have the most profound impact on the system's state.

Across all configurations tested, the network consistently returned to the local solution
\texttt{0001} after the kick was removed. This null result is itself informative:
single-pulse perturbations of the magnitudes and shapes tried were insufficient to escape
this particular local minimum, suggesting that future work should explore combinations
of parameters --- particularly amplitude, target neuron selection, and timing --- rather
than waveform shape alone.

% ─────────────────────────────────────────────────────────────────────────────
\section{Lessons Learned}
\label{sec:lessons}

\subsection{Understanding FPAA/Analog Hardware Limitations}

One of the most important realizations of the semester was how significantly physical
analog behavior can deviate from simulation. Transistor mismatch, floating-gate programming
precision, and temperature sensitivity caused results to vary between runs---the same weight
matrix and forcing inputs could produce completely different results. Running the same experiment under the same conditions at simply different times of the day would often yield different results.

In addition, the two FPAA
boards available in the lab did not behave identically. Even with the same compiled
configuration and programming sequence, the same circuit would produce measurably
different neuron voltages and convergence behavior on each board. Second, there also existed variability between individual I/O pins on the same board.
After discovering this, consistent usage of the same I/O pins was ensured.


\subsection{Communicating Frequently with Hari and Usman}

At the start of the semester, I found myself often trying to troubleshoot problems solely on my own, which led to lots of confusion which could've been easily cleared up by asking for help from Hari or Usman. For example, a significant amount of time was lost troubleshooting strange oscilloscope capture behavior that turned out to be caused by a typo in the FPAA board's
pinout documentation. Several of the document output pin mappings did not match the physical board, so the 
wrong signals were being read. This kind of error is nearly
impossible to diagnose from first principles---the circuit looks correct on paper, the
software compiles cleanly, and there is nothing in the experimental output that points
directly to a documentation error. Asking Hari and Usman resolved it in a single
conversation.

A related frustration was reproducibility: when attempting to replicate results that
Hari or Usman had previously obtained, different outcomes would appear despite using
theoretically identical configurations. This was consistent with the board-to-board and pin-to-pin variability described in
Section~\ref{sec:lessons}.1---their results may have been obtained on a different board under different conditions.  

\subsection{Documenting Everything}

A habit that took too long to develop was consistent documentation of experimental
results, regardless of whether they were satisfying. Early in the semester, oscilloscope
captures were only saved when the result seemed clearly meaningful or worth showing. Captures that were "failures" or did not show what was expected were often discarded without being recorded.

In hindsight, this was a mistake. Inconclusive results still carry information and can often point directly toward the next experiment worth trying. Several times during
the semester, a question arose about what had been tried before, and the answer was simply
unknown because no record existed. Going forward, every run---successful or not---should
be logged with the full parameter set, a CSV export of the oscilloscope data, and a brief
note on what was observed.

\subsection{Leveraging AI Tools Effectively}

AI tools proved to be very helpful in accelerating learning and automating tedious, repetitive tasks. The highest-impact use was
building a \href{https://gtvault-my.sharepoint.com/:u:/g/personal/emazor7_gatech_edu/IQDh167ysJgUTLz5orXGIl0TAVIHlGKhUOu2n0BPOJSZWZk?e=WR4lVW}{weight matrix generator} with Claude: manually converting a graph
topology into the correct $2N \times (2+N)$ conductance matrix---accounting for the
differential pair sign convention, conductance levels, budget constraints, and Scilab
output format---was error-prone by hand. Automating this eliminated an entire category of
repetitive mistakes and freed attention for analyzing results. {\small\color{darkgray}(\textit{Note: the link requires a Georgia Tech account; download the file and open it in a browser.})}

A second effective use was data visualization. The Analog Discovery exports
oscilloscope captures as CSV files, and feeding multiple CSVs into Gemini to generate
clean, consistently formatted plots proved far faster than manually processing them
in Excel or writing plotting scripts from scratch. This was particularly useful when
comparing neuron voltage traces across many kick signal configurations.

A third use was accelerating paper comprehension. Uploading research papers into a
Claude Project or a NotebookLM notebook allowed targeted questions and simple explanations of difficult or unfamiliar concepts. Finally, AI-assisted brainstorming proved useful for generating
candidate kick signal strategies---the parameter list in Section~\ref{sec:localvsglobal}
was developed with the help of Claude/Gemini.


\end{document}